\documentclass[10pt,letterpaper]{article}
\usepackage[margin=0.90in]{geometry}
\geometry{left=0.90in,bottom=0.90in}
\usepackage{enumitem}
\usepackage{tabularx}
\usepackage{hyperref}
\usepackage{titlesec}
\usepackage{lmodern}
\usepackage[T1]{fontenc}
\hypersetup{colorlinks=true,urlcolor=black,linkcolor=black}
\setlength{\parindent}{0pt}\setlength{\parskip}{0pt}
\titleformat{\section}{\large\bfseries}{}{0pt}{}[\vspace{-4pt}\rule{\textwidth}{0.4pt}]
\titlespacing*{\section}{0pt}{8pt}{5pt}
\setlist[itemize]{leftmargin=1.15em,itemsep=1.5pt,topsep=1pt,parsep=0pt,partopsep=0pt}
\newcommand{\heading}[2]{\vspace{2pt}\begin{tabularx}{\textwidth}{@{}Xr@{}}\textbf{#1} & \textit{#2}\end{tabularx}\vspace{-0pt}}
\newcommand{\headingwtitle}[3]{\vspace{2pt}\begin{tabularx}{\textwidth}{@{}Xr@{}}\textbf{#1} & #3\\[1pt] \textit{#2} & \end{tabularx}\vspace{-0pt}}

\begin{document}
\begin{center}
{\LARGE \textbf{Kebei Jiang}}\\[3pt]
Bellevue, WA \quad $\cdot$ \quad
\href{mailto:kebei.jiang13@gmail.com}{kebei.jiang13@gmail.com} \quad $\cdot$ \quad
(225) 200-4286 \quad $\cdot$ \quad
\href{https://www.linkedin.com/in/kebeijiang}{linkedin.com/in/kebeijiang}
\end{center}
\vspace{-3pt}

\section*{Profile}
Applied Scientist with 6+ years at Microsoft specializing in LLM/VLM foundation models, parameter-efficient fine-tuning, and agentic AI architectures. Experienced in controlled experimentation, model-behavior diagnostics, automated calibration, and systematic agent evaluation. Strong foundation in model training and optimization, with experience conducting empirical ML research and publishing peer-reviewed work.

\section*{Core Expertise}
\textbf{Foundation Models:} LLMs $\cdot$ VLMs $\cdot$ Parameter-Efficient Fine-Tuning $\cdot$ Agentic Systems

\textbf{Research \& Evaluation:} Controlled Experimentation $\cdot$ LLM-as-a-Judge Calibration $\cdot$ Systematic Agent Evaluation $\cdot$ Synthetic Data Generation $\cdot$ Evaluation Benchmarks $\cdot$ Statistical Analysis

\textbf{Model Training \& Optimization:} PyTorch $\cdot$ Transformers $\cdot$ LoRA $\cdot$ SFT $\cdot$ DPO $\cdot$ GRPO $\cdot$ Distributed Training $\cdot$ vLLM

\section*{Experience}
\headingwtitle{Microsoft}{Senior Applied Scientist}{Jan 2020 -- Sep 2026}
\begin{itemize}
\item Developed and deployed LLM and VLM solutions spanning agentic AI, RAG, knowledge mining, information extraction, prompt engineering, and model fine-tuning.
\item Designed evaluation methodologies and experiments for LLM-based agents and AI systems, including synthetic benchmarks, behavioral diagnostics, quantitative metrics, and regression evaluation.
\item Developed production ML/AI systems in Python, with experience in FastAPI, Docker, automated testing, production monitoring, and collaborative code reviews.
\item Built end-to-end ML and information-processing systems from data preparation and feature engineering through model development, evaluation, and deployment.
\item Led technical development across scientists, PMs, engineers, architects, and customer engineers, including technical design, implementation reviews, technical decisions, and customer engagements.
\end{itemize}

\headingwtitle{Viridien}{Applied Scientist}{Jan 2015 -- Jan 2020}
\begin{itemize}
\item Applied statistical and machine-learning methods to process and model large-scale (TB) seismic data.
\item Developed CNN-based solutions for seismic interpretation, including pixel-level identification of subsurface salt structures.
\end{itemize}

\section*{Selected Projects \& Research}

\heading{Specification-to-Calibration: Synthetic Supervision for LLM Judges}{2026}
\begin{itemize}
\item Designed and evaluated a novel algorithmic approach for automated LLM judge calibration by compiling natural-language specifications into coverage-controlled behavioral contrast suites.
\item Utilized feedback-conditioned and GEPA-based optimization to empirically demonstrate model adaptation, improving position-debiased accuracy by 3.58 pp on LLMBar benchmarks.
\end{itemize}

\heading{Period-Closing Skill --- Stateful, Multi-Turn Agent Evaluation}{2026}
\begin{itemize}
\item Developed an agentic skill for automating period-closing workflows involving due dates, calendars, worker assignments, and task dependencies.
\item Built an evaluation framework with 30+ reusable, parameterized scenario templates covering diverse failure modes; collaborated with the SMEs to validate scenario realism and coverage.
\item Created stateful, multi-turn test environments with pre-existing tasks, calendar constraints, and similarly named workers; simulated user responses to assess agent behavior across interaction turns.
\item Evaluated both final task outcomes and agent execution traces across repeated runs; measured task success rate (TSR), tool-call accuracy, latency, token cost, and consistency.
\end{itemize}

\heading{VLM Fine-Tuning --- Taxonomy Generalization \& Shortcut-Learning Diagnostic}{2025}
\begin{itemize}
\item Designed a hierarchical taxonomy spanning hundreds of fashion-item classes and aligned it with GPT-4o's understanding of retail products; generated $\sim$1M synthetic training examples 
\item Fine-tuned a Microsoft Phi-4-multimodal foundation model (5.6B) using PyTorch across multiple GPUs with LoRA, SFT, and DPO. Achieved $\sim$90\% accuracy against GPT-4o ground truths
\item Designed a rigorous, controlled diagnostic to analyze shortcut learning and taxonomy memorization, evaluating the limits of domain adaptation versus base model behavior.
\end{itemize}

\heading{Multimodal Retrieval \& Vision-Language Modeling}{2024}
\begin{itemize}
\item Fine-tuned a Phi-3 Vision model on public chart datasets to optimize complex visual question answering performance.
\item Evaluated cross-modal alignment and retrieval strategies using CLIP and SigLIP embeddings across distinct text, table, and image index spaces.
\item Researched advanced structural retrieval frameworks, implementing query decomposition, ranking fusion, and tree/graph-based architectures including RAPTOR and GraphRAG.
\end{itemize}

\section*{Advanced AI Infrastructure \& Systems Engineering}
\heading{CS336: Language Modeling from Scratch, Stanford}{2025}
\begin{itemize}
\item \textbf{Model \& Tokenizer Architecture:} Implemented all core components of a Transformer language model from scratch in PyTorch, including custom tokenization pipelines and optimization loops.
\item \textbf{Distributed Training \& Triton Kernels:} Built a memory-efficient, distributed version of the training code and optimized attention layers by writing custom Triton kernels for FlashAttention-2.
\item \textbf{Web-Scale Data Engineering:} Processed raw Common Crawl dumps into pretraining data by developing heuristic filtering and text deduplication pipelines.
\item \textbf{Empirical Scaling Laws:} Profiled layer execution traces and queried training APIs to mathematically fit scaling laws to project downstream performance.
\item \textbf{Reasoning \& Alignment Algorithms:} Programmed post-training alignment workflows, implementing SFT, DPO, and GRPO to train models to reason through math problems.
\end{itemize}

\section*{Publications}
\begin{itemize}
\item Guang Yang, Kebei Jiang, David P. Lundquist, Prerna Singh. \textit{Specification-to-Calibration: Synthetic Supervision for LLM Judges.} Manuscript in Preparation.
\item Moschoula Pternea, Prerna Singh, Abir Chakraborty, Yagna Oruganti, Mirco Milletari, Sayli Bapat, Kebei Jiang. \textit{The RL/LLM Taxonomy Tree: Reviewing Synergies Between Reinforcement Learning and Large Language Models.} Journal of Artificial Intelligence Research, 2024. DOI: 10.1613/jair.1.15960.
\item Yu Zeng, Kebei Jiang, Jie Chen. \textit{Automatic Seismic Salt Interpretation with Deep Convolutional Neural Networks.} Proceedings of the 2019 3rd International Conference on Information System \& Data Mining, pp. 16--20.
\end{itemize}

\section*{Patent}
\begin{itemize}
\item Yinheng Li, Kebei Jiang. \textit{Concurrent Labeling of Sequences of Words and Individual Words.} U.S. patent application, MS\# 411939-US-NP, 2022.
\end{itemize}

\section*{Education}
\headingwtitle{Louisiana State University, Baton Rouge, LA}{PhD, Physics}{2007 -- 2014}

\section*{Personal Projects}

\heading{OptiBuild --- Domain-Agnostic Optimization Agent | Personal Project}{2026}
\begin{itemize}
\item Engineered a domain-agnostic multi-agent system that translates natural-language constraints into structured Operations Research models, enabling optimized component selection across arbitrary domains (e.g., PC hardware, drone components).
\item Developed an orchestration pipeline utilizing Gemini and CP-SAT with TOPSIS-based multi-objective ranking to evaluate intent fidelity and systematically solve for Pareto-optimal configurations.
\item Designed a staged modelization and repair loop utilizing Pydantic contracts, deterministic validation protocols, and structured execution traces to handle infeasible models or missing data dynamically.
\item Implemented rigorous automated testing and security controls, incorporating a 127-case offline regression suite, PII redaction layers, and strict safety gates for execution-level LLM queries and formulas.
\end{itemize}


\end{document}
